% Transcription — fonds Alexandre Grothendieck, shelfmark 21, batch 1 (pages 1-20).
% Established from the facsimile archives/batches/21/batch-01.pdf (a mirror of
% https://grothendieck.umontpellier.fr/21.pdf), per the skill
% transcribe-grothendieck. The batch file carries no cover sheet: PDF page N is
% archivists' page N. Checked against the pencilled numbers bottom-left — a "1"
% on the first leaf, then "8", "10", "11", "13", "14", "15", "17", "18", "19",
% "20" on the corresponding leaves, confirm the alignment end to end.
%
% Pass: Opus 5 (claude-opus-5), 2026-09-19 — first pass, unchecked against the pages by a human.
%
% What the batch is. Not a manuscript: two things wrapped in kraft.
%
%   Pages 3 to 13 are eleven leaves of a carbon typescript of his own, headed
%   « 2. Généralités sur certains foncteurs de modules. », running from
%   Proposition 2.1 to Remarques 2.17 without a break and unpaginated by him.
%   The body is typed and legible; what carries the interest of the piece is
%   the correction layer, which is of two distinct kinds and is transcribed as
%   two distinct things: deletions made ON THE MACHINE, by overstriking a word
%   with x's (\struck{}), and additions made LATER BY HAND, in ink, between the
%   lines or in the margin (\add{}, \marginal{}). A third kind — a word retyped
%   above the line in the same typewriter face, which is neither — is recorded
%   in a \note{} where it occurs.
%
%   Pages 15 to 20 are six leaves in his hand, a separate run headed
%   « Propriétés (P1) et (P2) pour des faisceaux d'ensembles », numbered by him
%   Définition 1, Définition 2, Proposition 3, Corollaire 4, Lemme 5. This is a
%   fast draft: the statements and the displayed formulas come through, the
%   connective prose very often does not, and the apparatus says so rather than
%   guessing. Pages 17 and 19 in particular carry long cancelled and rewritten
%   stretches that this pass could not read.
%
% Pages skipped, and why — the gap in the numbering is the record:
%   1   kraft wrapper, blank but for « Notes pour Dieud. §9 » in his hand,
%       top right. The user's brief gave this as « §3 »; the leaf reads 9, and
%       the digit is clean enough to say so. No mathematics on the sheet.
%   2   verso of leaf 3 — nothing but the show-through of the typescript.
%   12  verso of leaf 11 — likewise, show-through only.
%   14  second kraft wrapper, blank but for six lines in his hand, struck
%       through with three long diagonals: « Foncteurs \uncertain{élémentaires}
%       / sur les espaces topologiques, / en particulier noethériens, /
%       profondeur >= 2 / pour faisceaux / d'ensembles ». It names the run of
%       pages 15-20. Per the brief a wrapper takes no \page and no \note, so it
%       is recorded here and nowhere else; the second section heading is taken
%       from page 15's own first line instead.
%
% Nothing in the batch is dated and no date appears on any leaf. The
% inventory's « [vers 1960-1967] » is carried into \dating{} verbatim and is
% the archivists', not his.
%
% The typescript's own numbering (2.1 to 2.17) is his draft's and is
% transcribed as it stands. It is NOT reconciled with printed EGA, and printed
% EGA was not opened for this pass.
%
% Single overstruck characters — a letter typed, x-ed out and retyped — run
% through the whole typescript and are typing slips, not deletions. They are
% not recorded. Only overstrikes that cancel a word or more appear as
% \struck{}.
\documentclass[11pt,a4paper]{article}
% Le préambule commun à toutes les transcriptions du fonds Grothendieck.
%
% Il tient en une idée : **l'incertitude se compose, elle ne se résout pas.**
% Une transcription qui lisse un mot douteux a détruit la seule chose pour
% laquelle on la faisait — un lecteur doit pouvoir savoir, à chaque endroit,
% ce qui était sur le feuillet et ce qui a été ajouté. Les six macros qui
% suivent sont donc l'appareil critique, et non de la décoration.
%
% Les mêmes commandes sont reconnues par `scripts/render.mjs`, qui produit la
% vue de lecture du site. Une macro ajoutée ici doit l'être aussi là-bas,
% faute de quoi le rendu échouera bruyamment — ce qui est voulu : un rendu
% silencieusement faux est pire qu'un rendu absent.

\ProvidesPackage{grothendieck}[2026/08/08 Transcriptions du fonds Grothendieck]

\usepackage{amsmath,amssymb,amsthm}
\usepackage{mathrsfs}
\usepackage{stmaryrd}
% Les diagrammes commutatifs sont partout dans ce fonds : c'est la langue
% même de la géométrie algébrique de Grothendieck, et une transcription qui
% les rendrait en prose ne transcrirait rien.
\usepackage{tikz-cd}
% Français par défaut — c'est la langue des feuillets — mais l'anglais est
% chargé aussi : la traduction et le résumé sont des documents anglais, et la
% typographie française (espaces avant les deux-points) y serait une faute.
% Ces documents basculent par \selectlanguage{english} en tête de corps.
\usepackage[english,french]{babel}
\usepackage{fontspec}
\usepackage{geometry}
\geometry{a4paper,margin=25mm,marginparwidth=28mm}
\usepackage{marginnote}
\usepackage{xcolor}
% ulem, not soul. `soul`'s \st and \ul are built on a character-by-character
% scan that XeLaTeX's font machinery breaks: the document fails with an
% undefined control sequence rather than degrading. ulem does the same two jobs
% and is Unicode-engine safe. `normalem` stops it hijacking \emph, which the
% transcriptions use for Grothendieck's own emphasis.
\usepackage[normalem]{ulem}
\usepackage{hyperref}
\hypersetup{colorlinks=true,linkcolor=black,urlcolor=black,pdfborder={0 0 0}}

% --- Métadonnées du lot -----------------------------------------------------
% Reprises telles quelles par le script de rendu, qui les affiche en tête de
% la vue de lecture. Elles fixent ce que le fichier prétend couvrir : sans
% elles, une transcription flotte sans point d'attache dans un fonds de seize
% mille feuillets.

\newcommand{\folder}[1]{\gdef\grothFolder{#1}}
\newcommand{\batch}[1]{\gdef\grothBatch{#1}}
\newcommand{\leaves}[2]{\gdef\grothFirst{#1}\gdef\grothLast{#2}}
\newcommand{\foldertitle}[1]{\gdef\grothFoldertitle{#1}}
\gdef\grothFolder{?}\gdef\grothBatch{?}\gdef\grothFirst{?}\gdef\grothLast{?}\gdef\grothFoldertitle{}

% --- Le résumé --------------------------------------------------------------
%
% Il ouvre la lecture modernisée, sous le titre « Résumé », et il est composé
% en petit. Ce n'est pas un ornement : c'est le seul passage du document qui
% ne soit pas des mathématiques, et un lecteur doit pouvoir le reconnaître
% avant de l'avoir lu — puis le sauter s'il connaît déjà le sujet, ou s'y
% arrêter s'il ne le connaît pas. Le filet qui le suit marque la reprise du
% corps.

\newenvironment{resume}
  {\small\setlength{\parskip}{0.45em}}
  {\par\medskip\hrule\medskip}

% --- Mots-clés ---------------------------------------------------------------
%
% En anglais, en clôture du résumé de la lecture modernisée : le vocabulaire
% moderne sous lequel chercher ce que les feuillets construisent. Le manifeste
% du site (`npm run manifest`) les extrait pour étiqueter la cote entière —
% cette ligne est la source unique des tags, il n'y a pas de fichier à côté.
\newcommand{\keywords}[1]{\par\smallskip\noindent
  {\footnotesize\textsc{Keywords} --- \textit{#1}}\par}

% --- L'appareil critique ----------------------------------------------------

% Le numéro de feuillet, en marge. C'est l'ancre qui permet de régler un
% désaccord sur une formule en regardant *un* feuillet plutôt qu'une chemise
% de six cents. Le site s'en sert aussi pour tourner les pages du fac-similé
% quand on fait défiler la transcription.
\newcommand{\leaf}[1]{%
  \marginnote{\footnotesize\color{gray!70}\textsf{#1}}%
  \label{leaf:#1}\ignorespaces}

% Illisible : on ne devine pas. Un mot inventé qui se lit comme les autres est
% le pire résultat possible.
\newcommand{\ill}{\textcolor{red!60!black}{[\,\ldots\,]}}

% Lecture douteuse : le mot est donné, et signalé comme incertain.
\newcommand{\uncertain}[1]{\uline{#1}}

% Ajout de l'éditeur — une lettre manquante, un mot sous-entendu.
\newcommand{\add}[1]{\textcolor{blue!50!black}{[#1]}}

% Biffé par Grothendieck lui-même. Ce qu'il a rayé fait partie du document :
% une rature dit souvent où la pensée a changé de direction.
\newcommand{\struck}[1]{\sout{#1}}

% Note du transcripteur, sur l'état matériel du feuillet.
\newcommand{\note}[1]{\par\smallskip{\small\color{gray!60!black}\textit{[#1]}}\par\smallskip}

% Note marginale de Grothendieck, distincte du corps du texte.
\newcommand{\marginal}[1]{\par\smallskip\noindent\hspace{1em}%
  {\small\color{gray!60!black}#1}\par\smallskip}

% --- Titre ------------------------------------------------------------------

\AtBeginDocument{%
  \begin{center}
    {\large\bfseries Fonds Alexandre Grothendieck --- cote n\textsuperscript{o}~\grothFolder}\\[2pt]
    {\itshape\grothFoldertitle}\\[4pt]
    {\small feuillets \grothFirst--\grothLast\ (lot \grothBatch)}\\[2pt]
    {\footnotesize\color{gray!60!black}Transcription établie d'après le fac-similé numérisé
     par l'Université de Montpellier.\\
     Les lectures incertaines sont soulignées, les illisibles notées [\,\ldots\,],
     les ajouts entre crochets.}
  \end{center}
  \vspace{1em}}

\endinput


\folder{21}
\batch{1}
\pages{1}{20}
\dating{[vers 1960-1967]}
\watermark{Édition de démonstration}
\foldertitle{[Chapitre] 0 III : tapuscrit et copies de tapuscrit annotés (s.d.), notes manuscrites (s.d.).}

\begin{document}

\section{2. Généralités sur certains foncteurs de modules}

\page{3}
2. Généralités sur certains foncteurs de modules.

Le présent numéro pourrait passer également au \struck{\ill{}} Chap $\mathrm{O}_{III}$, qui contiendrait aussi certains des sorites du par.7.

\emph{Proposition 2.1.} Soient $\underline{C}$ une catégorie abélienne, $\underline{G}$ une partie de $\mathrm{Ob}\,\underline{C}$, telle que tout objet de $\underline{C}$ admette une suite de composition finie dont les \struck{quotients} facteurs \note{« facteurs » est retapé au-dessus de la ligne, dans la même frappe} sont isomorphes à des quotients d'objets $M \in \underline{G}$. \struck{\ill{}} Soit $\underline{C}'$ une seconde catégorie abélienne.

(i) Soit $T : \underline{C} \to \underline{C}'$ un foncteur \struck{\ill{}} exact à droite, pour que $T$ soit nul, il f et s qu'il soit nul sur les $M \in \underline{G}$.

(ii) Soit $(T^i)$ un foncteur cohomologique de $\underline{C}$ dans $\underline{C}'$, tel que $T^i = 0$ pour $i > i_o$, et soit $n$ un entier. Pour qu'on ait $T^i = 0$ pour $i > n$, il f. et s. que pour tout $M \in \underline{G}$, on ait $T^i(M) = 0$ pour $i > n$.

Dém. (i) est trivial, pour (ii) on procède par récurrence descendante \note{« descendante » est retapé au-dessus de la ligne, avec un signe d'insertion à l'encre} sur $n$, c'est trivial pour $n \geqslant i_o$, et pour $n$ quelconque, supposant que c'est fait pour $n+1$, on note que si les $T^i(M)$ ($M \in \underline{G}$, $i \geqslant 1+n$) sont nuls, alors par l'hy de réc on a $T^i = 0$ pour $i > n+1$, donc $T^{n+1}$ est exact à droite, et d'après (i) il est nul, cqfd.

\emph{Corollaire 2.2.} \note{le premier chiffre est surchargé à la machine : un 3 et un 2 superposés} Sous les conditions de 2.1. soit $(T^i)$ un foncteur cohomologique contravariant de $\underline{C}$ dans $\underline{C}'$, tel que $T^i = 0$ pour $i < i_o$, et soit $n$ un entier. Pour qu'on ait $T^i = 0$ pour $i < n$, il f et s que pour tout $M \in \underline{G}$, on ait $T^i(M) = 0$ pour $i < n$.

Il suffit d'appliquer 2.1. (ii) à $\underline{C}$ et $\underline{C}'^{o}$. Pour qu'il n'y ait pas de jaloux, il faudrait aussi répéter (i). Dans le présent paragraphe, ce sera sous cette forme que nous utiliserons 2.1. (mais il aurait été canularesque d'énoncer 2.1. avec des foncteurs contravariants).

\emph{Corollaire 2.3.} Soient $\underline{C}$, $\underline{C}'$ des catégories abéliennes, $\underline{G}$ une partie de $\mathrm{Ob}\,\underline{C}$, on suppose que les objets de $\underline{C}$ sont noethériens (i.e. satisfont la condition des chaines ascendantes pour les sous-objets), et que $\underline{G}$ est un systême de générateurs de $\underline{C}$ (i.e. que pour tout $F \neq 0$ dans $\underline{C}$, il existe un $M \in \underline{C}$ et un homomorphisme non nul $M \to F$). Alors les conclusions de 2.1. (i)(ii) et 2.2. (i)(ii) sont valables.

Il suffit de vérifier les hypothèses de 2.1. \add{sur $\underline{G}$}, pour ceci soit $F \in \mathrm{Ob}\,\underline{C}$, \struck{\ill{}} soit $F'$ un sous-objet maximal de $F$ admettant une suite de composition finie comme dite dans 2.1. (un tel $F'$ existe, car $0 < F$ satisfait à l'hypothèse de 2.1., et on applique l'hyp noeth sur $F$). Je dis que $F' = F$, en

\page{4}
effet sinon on aurait $F/F' \neq 0$, donc il existerait un homomorphisme non nul $M \to F/F'$ ($M \in \underline{G}$), et l'image inverse \add{dans $F$} de son image \struck{dans} contredirait l'hypothèse maximale pour $F'$.

\emph{Exemple 2.4.} Soient $A$ un anneau comm noeth, $I$ un idéal de $A$, $X$ le spectre de $A$ et $Y$ la partie fermée de $X$ définie par $I$. Soit $\underline{C}$ la catégorie des $A$-modules \add{de type fini} dont tout élément est annulé par une puissance de $I$, ou la catégorie équivalente des Modules \struck{quasi-}cohérents sur $X$ dont le support est contenu dans $Y$. Les objets de $\underline{C}$ sont noethériens. Si $M \in \mathrm{Ob}\,\underline{C}$, alors on voit tout de suite que $M$ est un générateur de $\underline{C}$ sss $\mathrm{supp}\, M = Y$.

C'est nécessaire, car si on avait $\mathrm{supp}\, M \neq Y$, il existerait une composante irréductible $Y_i$ de $Y$ non contenue dans $\mathrm{supp}\, M$, et prenant $F = \mathcal{O}_{Y_i}$, avec $Y_i$ muni de la structure réduite induite, on trouve que $\mathrm{Hom}(M,F) = 0$, $F \neq 0$.

C'est suffisant, car si $F$ est quasi-cohérent, de support $\subset Y$, on a $\mathrm{Ass}\,\mathcal{H}om(M,F) = \mathrm{supp}\, M \cap \mathrm{Ass}\, F = \mathrm{Ass}\, F$, donc $\mathrm{Hom}(M,F) = 0$ implique $\mathrm{Ass}\, F = \emptyset$ d'où $F = 0$. (Donc $M$ est un générateur également dans la catégorie de tous les modules quasi-cohérents $F$ de support $\subset Y$). \struck{Par suite, \ill{} pour tester la nullité de} Le même argument montre que si on a une famille $(M_i)_{i \in I}$ d'objets de $\underline{C}$, alors c'est un système de générateurs de $\underline{C}$ sss la réunion des $\mathrm{supp}\, M_i$ est $Y$ (commencer par ça). Pour une famille satisfaisant cette condition, l'hypothèse de 2.1. est vérifiée, et par suite, pour vérifier la nullité de foncteurs comme dans 2.1. ou 2.2., il suffit de tester par les $M_i$, \struck{ou de tester \ill{}} en particulier il suffit de tester par un $M$ de support $Y$.

\emph{Exemple 2.5.} On peut généraliser 2.4. en prenant un schéma \struck{$X$ muni d'un} quasi-compact \note{« quasi-compact » est retapé au-dessus du passage surfrappé} noethérien $X$, muni d'un faisceau ample noté $\mathcal{O}_X(1)$, \struck{Soit $(M_i)$} et d'une partie fermée $Y$, $\underline{C}$ désignant encore la catégorie des faisceaux cohérents sur $X$ à support $\subset Y$. Les objets de $\underline{C}$ sont encore noethériens, et si $(M_i)_{i \in I}$ est une famille d'objets de $\underline{C}$, telle que \struck{$\mathrm{supp}\, M_i \ill{}$}
\[
  \bigcup_i \mathrm{supp}\, M_i \;=\; Y ,
\]
(il aurait fallu déployer la formule également dans 2.4.), alors les $M_i(-n)$ \note{le signe de l'exposant est surchargé à la machine}, avec $n > N$ ($N$ entier donné), forment un système de générateurs de $\underline{C}$ (et même de la catégorie des Modules qu coh de support $Y$). En effet si $F$ est un Module qu coh de support $Y$, et $F \neq 0$, alors \struck{\ill{}} on a vu dans 2.4. que le faisceau $\mathcal{H}om(M_i,F)$ est non nul pour $i$ convenable (se placer dans un ouvert affine où $F$ est $\neq 0$), donc il existe des sections non nulles de $\mathcal{H}om(M_i,F)(n)$ pour $n$ grand, donc des homomorphismes non nuls de $M_i(-n)$ dans $F$ pour $n$ grand. -- On peut donc utiliser les $M_i(-n)$, $n \geqslant N$,

\page{5}
pour tester la nullité de foncteurs comme dans 2.1. et 2.2.

Soient maintenant $A$ un anneau (non nécessairement commutatif), $\underline{C}$ une catégorie abélienne de $A$-modules à gauche \note{la page dit bien « catégorie abélienne de $A$-modules », et non « la catégorie »}, par exemple n'importe quelle sous-catégorie pleine \note{« pleine » est retapé au-dessus de la ligne} de la catégorie des $A$-modules, stable par noyaux et conoyaux, $\underline{C}'$ une catégorie abélienne. Nous allons construire, pour tout objet $P$ de $\underline{C}'$ muni d'une structure de $A$-module à droite, i.e. muni d'un homomorphisme d'anneaux
\[
  A^{o} \longrightarrow \mathrm{End}(P) ,
\]
et tout $M \in \mathrm{Ob}\,\underline{C}$, un objet de $\underline{C}'$ noté
\[
  P \otimes_A M ,
\]
défini par la propriété universelle
\[
  \mathrm{Hom}_{\underline{C}'}(P \otimes_A M,\, Q) \;=\; \mathrm{Hom}_A(M,\, \mathrm{Hom}_{\underline{C}'}(P,Q)) ,
\]
le second membre ayant un sens grâce à la structure de $A$-module à gauche de $\mathrm{Hom}_{\underline{C}'}(P,Q)$ déduite de celle de $A$-module à droite de $P$. Cet objet n'existe pas toujours, sa construction \struck{définition} \note{« construction » est retapé au-dessus du mot surfrappé} utilisera le fait que \struck{\ill{}} $\underline{C}$ contient $A$ et que les objets de $\underline{C}$ sont des modules de présentation finie sur $A$.

\marginal{déployer}

Notons d'abord que si $M = A_g$, $A$ considéré comme $A$-module à gauche, et regardé comme \struck{l'unique} objet de $\underline{C}'$, alors $P \otimes_A A_g$ existe et est canoniquement isomorphe à $P$.

\page{6}
D'après les sorites généraux sur les foncteurs représentables, $P \otimes_A M$ dépend fonctoriellement de $(M,P)$ \struck{et de $P$,} \struck{i.e. \ill{} il} pour les \add{couple d'}arguments pour lesquels il est défini.

De plus, on voit tout de suite que si on a une suite exacte
\[
  M'' \longrightarrow M' \longrightarrow M \longrightarrow 0
\]
de $A$-modules à gauche, telle que $P \otimes_A M'$ et $P \otimes_A M''$ existent, alors $P \otimes_A M$ existe et \struck{on a une} \add{la} suite \struck{exacte}
\[
  P \otimes_A M'' \longrightarrow P \otimes_A M' \longrightarrow P \otimes_A M \longrightarrow 0
\]
\add{est exacte} (en particulier, $P \otimes_A M$ est foncteur exact à droite en $M$, sous réserve d'existence). \note{« foncteur exact à droite en $M$ » est souligné d'un trait ondulé à l'encre}

On a le fait analogue pour des suites exactes en $P$. D'autre part, $P \otimes_A A_g$ existe toujours et on a un isomorphisme canonique
\[
  P \otimes_A A_g \;\simeq\; P .
\]
\marginal{et plus gén. des $\varinjlim$ \ill{}} \note{les deux dernières phrases sont cernées d'un trait à l'encre, la note marginale s'y rattache}

\struck{(Pour} (Tous ces faits sont triviaux sur les définitions, en interprétant en termes de foncteurs en $Q$). \struck{Cela prouve\ill{}} Plus généralement, le foncteur $P \otimes_A M$ en $M$, ou en $P$, commute aux limites inductives quelconques (pas nécessairement filtrantes), au sesn \note{sic} suivant: si les $P \otimes_A M_i$ existent, \struck{et $\varinjlim$} alors $P \otimes_A \varinjlim M_i$ existe sss $\varinjlim P \otimes_A M_i$ existe, et les deux sont can isomorphes.

Ceci s'applique en particulier aux limites inductives sur un ensemble discret, i.e. aux sommes directes.

Combinant ces remarques, on voit que $P \otimes_A M$ existe \struck{toujours} dès que $M$ est de présentation finie (car il existe pour $M = A_g^{(I)}$, $I$ ensemble fini, donc pour des conoyaux d'homomorphismes $A_g^{(I)} \to A_g^{(J)}$, $I$ et $J$ finis), et sans restriction sur $M$ si dans $\underline{C}'$ les sommes directes quelconques existent, (même argument, sans restriction de finitude sur $I$, $J$). De plus, on a un procédé de construction de $P \otimes_A M$ par « générateurs et relations » en termes de « générateurs et relations » pour $M$, i.e. en termes \struck{d'un} de la matrice définissant une présentation \struck{finie} de $M$ (cette matrice définit une matrice sur $I \times J$, définissant une application $P^{(I)} \to P^{(J)}$, dont le conoyau est le produit tensoriel cherché). Tout ceci devrait figurer dans un numéro à part \add{au Chap 0}, sans trainer l'accent $'$ sur la catégorie abélienne. Du point de vue notations, si $\underline{C}'$ est \struck{la catégorie} l'opposée \note{« l'opposée » est retapé au-dessus du passage surfrappé} d'une catégorie abélienne $\underline{C}_1$, \struck{on note Hom} les $A$-$\underline{C}'$-modules à droite sont les $A$-$\underline{C}_1$-modules à gauche, et le produit tensoriel dans $\underline{C}'$ se note (en termes de $\underline{C}_1$)
\[
  \mathrm{Hom}_A(M,P)
\]
défini directement par \struck{$\mathrm{Hom}_Q(\mathrm{Hom}(M,P))$}

\page{7}
la propriété universelle
\[
  \mathrm{Hom}_{\underline{C}_1}(Q,\, \mathrm{Hom}_A(M,P)) \;=\; \mathrm{Hom}_A(M,\, \mathrm{Hom}_{\underline{C}_1}(Q,P)) .
\]
\note{les deux membres de cette formule sont les seuls que le feuillet porte en clair ; la ligne qui la précède est reconstituée du contexte}

Il y aurait donc lieu de développer les deux formalismes en parallèle, l'un portant sur des produits tensoriels, l'autre sur des Hom, \struck{sur les} modules habituels (il faut avoir vu bien entendu qu'on retrouve ces \struck{derniers lorsque} \add{\ill{} habituel, à} \note{sous la surfrappe, une ligne à l'encre dont seul « habituel » se lit} $\underline{C}'$ resp. $\underline{C}_1$ est la catégorie (Ab), i.e. comme cas particulier de produits tensoriels avec des $A$-(Ab)-modules et des $A$-(Ab)$^{o}$-modules).

Soit donc $\underline{C}$ la catégorie des $A$-modules à gauche de présentation finie, et soit $\underline{C}'$ une catégorie abélienne. Pour tout $A$-$\underline{C}'$-module à droite $P$ on a donc un foncteur additif $M \rightsquigarrow P \otimes_A M$ de $\underline{C}$ dans $\underline{C}'$, et un homomorphisme canonique, fonctoriel en $M$ :
\[
  P \otimes_A M \longrightarrow \ill{}
\]
\note{le second membre de ce diagramme n'a pas été lu}

\emph{Proposition 2.6.} Soit $T$ un foncteur additif de $\underline{C}$ dans $\underline{C}'$, et soit $P = T(A)$, muni de sa structure naturelle de $A$-module à droite. L'homomorphisme canonique $P \otimes_A M \to T(M)$ est un isomorphisme si et seulement si $T$ est exact à droite.

\note{la page 7 n'a pu être lue de bout en bout ; le texte ci-dessus en donne les énoncés, les phrases de liaison manquent}

\page{8}
contravariant $T$ de $\underline{C}$ dans $\underline{C}'$, i.e. d'un foncteur
\[
  T : \underline{C}^{o} \longrightarrow \underline{C}' ,
\]
(ce cas se ramenant au précédent en remplaçant $\underline{C}'$ par $\underline{C}'^{o}$). On trouve maintenant un homomorphisme canonique, fonctoriel en $M$ :
\[
  T(M) \longrightarrow \mathrm{Hom}_A(M, T(A)) ,
\]
et il faudra expliciter les énoncés duals de 2.6. et 2.7., en particulier l'homomorphisme précédent est un isomorphisme sss $T$ est exact à gauche (i.e. transforme une suite exacte $M'' \to M' \to M \to 0$ dans $\underline{C}$ en une suite exacte $0 \to T(M) \to T(M') \to T(M'')$) : on obtient une équivalence entre la catégorie des foncteurs contravariants exacts à g de $\underline{C}$ dans $\underline{C}'$ et la catégorie \add{$A$}$(\underline{C}')$. Ces résultats s'étendent \add{comme dans \uncertain{2.4.}} au cas où $\underline{C}$ est la catégorie de tous les $A$-modules, sans hypothèse noethérienne sur $A$, pourvu que dans $\underline{C}'$ les produits quelconques existent.

Nous allons généraliser les considérations précédentes à la catégorie $\underline{C}$ de l'exemple 2.4. Soit pour tout $n \geqslant 0$, $\underline{C}^{n}$ la sous-catégorie pleine de $\underline{C}$ formée des modules de type fini sur $\underline{C}$ \note{la page porte « sur $\underline{C}$ »} annulés par $I^{n+1}$, catégorie équivalente à la catégorie des modules de type fini sur $A_n = A/I^{n+1}$. Alors $\underline{C}_n \to \underline{C}$ est un foncteur exact, et $\underline{C}$ est réunion filtrante de ses sous-catégories $\underline{C}_n$. Alors un foncteur \add{additif}
\[
  T : \underline{C} \longrightarrow \underline{C}'
\]
est défini par la suite des foncteurs restriction
\[
  T_n : \underline{C}_n \longrightarrow \underline{C}' ,
\]
et en vertu de ce qui précède, pour chaque $T_n$ on a un homomorphisme fonctoriel
\[
  T_n(A_n) \otimes_{A_n} M \longrightarrow T_n(M) = T(M)
\]
qu'on peut aussi écrire
\[
  \struck{\ill{}} \qquad T(A_n) \otimes_A M \longrightarrow T_n(M) \quad \add{(n \geqslant n_o)} .
\]
Pour $M$ fixé dans un $\underline{C}_{n_o}$, on obtient ainsi un système projectif \struck{\ill{}} $(T(A_n) \otimes_A M)_{n \geqslant n_o}$ d'objets de $\underline{C}'$, et un homomorphisme de ce dernier dans le système projectif constant défini par \struck{\ill{}} l'objet $T(M)$. Sous réserve de l'existence de $\varprojlim T(A_n) \otimes \struck{\ill{}} M$, on trouve donc un homomorphisme can
\[
  \varprojlim T(A_n) \otimes_A M \longrightarrow T(M)
\]
d'où par composition un homomorphisme canonique
\[
  \left(\varprojlim T(A_n)\right) \otimes_A M \longrightarrow T(M) ,
\]

\page{9}
et dans ces derniers homomorphismes, il n'est plus nécessaire de faire mention de l'entier $n_o$ (i.e. ils ne dépendent pas du choix de l'entier $n_o$ tel que $M \in \mathrm{Ob}\,\underline{C}_{n_o}$). Pour notre objet, nous aurons plutot à considérer la situation duale d'un foncteur contravariant sur $\underline{C}$ :
\[
  T : \underline{C}^{o} \longrightarrow \underline{C}' ,
\]
donnant lieu à des homomorphismes
\[
  (\%) \qquad T_n(M) \longrightarrow \mathrm{Hom}_A(M, T_n(A_n)) \qquad (n \geqslant n_o) ,
\]
satisfaisant les propriétés de transitivité qu'il faut pour pouvoir passer à la limite (sous réserve d'existence des $\varinjlim$ en question)
\[
  T(M) \longrightarrow \varinjlim_n \mathrm{Hom}_A(M, T(A_n)) \longrightarrow \mathrm{Hom}_A\!\left(M, \varinjlim_n T(A_n)\right)
\]
Pour que $T$ soit exact à gauche, il f et s que les $T_n$ le soient, ou encore que les homomorphismes $(\%)$ soient des isomorphismes. Si dans $\underline{C}'$ les limites inductives \add{filtrantes} \note{« filtrantes » est retapé au-dessus de la ligne, avec un caret} dénombrables existent et sont "exactes", à fortiori si l'axiome (Ab 5) est vérifié, alors \struck{l'exactitude à gauche de $T$ implique} \struck{que donc que l'homomorphisme} on voit facilement que le foncteur $P \rightsquigarrow \mathrm{Hom}_A(M,P)$ de \struck{$\underline{C}$ dans $\underline{C}$} $A(\underline{C}')$ dans $\underline{C}'$ commute aux limites inductives \add{(filtr)} dénombrables (pour tout $M$ de présentation finie sur $A$). Donc \struck{sous} si \add{de plus} $T$ est exact à gauche, l'homomorphisme composé
\[
  T(M) \longrightarrow \mathrm{Hom}_A\!\left(M, \varinjlim T(A_n)\right)
\]
est un isomorphisme, la réciproque étant triviale. Il y a lieu de résumer ce résultat dans une

\emph{Proposition 2.9.}

\note{l'énoncé de 2.9. n'est pas écrit : le feuillet porte le seul intitulé, puis reprend}

généralisant 2.8. sous les conditions dites, et en déduire le

\emph{Corollaire 2.10.} \note{le second chiffre est surchargé} $\underline{C}$, $\underline{C}'$ étant comme dans 2.8., la catégorie des

\page{10}
\note{la page s'ouvre en cours de phrase ; le début en est perdu avec la fin de la page 9}
morphisme de transition $H_n \to H_m$ induit un isomorphisme de $H_n$ sur le sous-$A$-$\underline{C}'$-module $\mathrm{Hom}_A(A_n, H_m)$ des "éléments annulés par $I^{n+1}$". Bien entendu, sous l'hypothèse sur $\underline{C}'$ faite dans 2.9. et 2.10. on voit aussitôt directement que la catégorie de systêmes inductifs qu'on vient de décrire, est équivalente à \struck{\ill{}} la catégorie de $A$-$\underline{C}'$-modules mentionnée dans 2.10., \add{qui est donc un cas particulier de 2.11.}

\emph{Corollaire 2.12.} Soient $A$ un anneau comm noeth. muni d'un idéal $I$, $\underline{C}$ la catégorie des $A$-modules de type fini de support $\subset V(I)$ i.e. annulés par une puissance de $I$, $\underline{C}'$ une catégorie abélienne à limites inductives filtrantes dénombrables exactes, $(T^i)$ un foncteur cohomologique contravariant de $\underline{C}$ dans $\underline{C}'$, tel que $T^i = 0$ pour $i < i_o$, $n$ un entier, $(M_j)$ une famille de $A$-modules de type fini tel que la réunion des supports des $M_j$ soit égale à $Y = V(I)$. \struck{\ill{}} Posons pour tout entier $i$:
\[
  H^i = \varinjlim_m T^i(A_m) \qquad (\text{où } A_m = A/I^{m+1}) .
\]
Alors les conditions suivantes sont équivalentes:
\begin{itemize}
  \item (i) On a $T^i = 0$ pour $i < n$.
  \item (ii) On a $H^i = 0$ pour $i < n$
  \item (iii) On a $T^i(M_j) = 0$ pour $i < n$ et tout $j$.
\end{itemize}

Dém. L'équivalence de (i) et (iii) est mise pour mémoire, et a été vue dans 2.4. On a (i) $\Rightarrow$ (ii), pour l'implication inverse on raisonne par récurrence sur $n$, c'est trivial pour $n \leqslant i_o$, pour $n$ quelconque et supposant que c'est fait pour les $n' < n$, il reste à prouver que $T^{n-1} = 0$, or $T^{n-1}$ est exact à gauche, donc en vertu de 2.9. est isomorphe à $\mathrm{Hom}_A(M, H^{n-1})$, \add{donc \uncertain{nul},} \note{lecture douteuse : le mot ajouté à l'encre peut se lire « nul » ou « une » ; c'est « nul » que l'argument demande} ce qui achève la démonstration. De plus, on voit:

\emph{Corollaire 2.13.} Si les conditions équivalentes (i) à (iii) de 2.12. sont vérifiées, on a un isomorphisme fonctoriel
\[
  T^{n}(M) = \mathrm{Hom}_A(M, H^{n})
\]
\note{l'exposant est surchargé à la machine aux deux occurrences ; « $n$ » est ce que porte le dernier état}

Il serait préférable d'inclure 2.13. dans l'énoncé de 2.12. et de démontrer par récurrence (ii) $\Rightarrow$ (i) \emph{plus} l'énoncé final 2.13., pour éviter des redites.

\page{11}
\emph{Corollaire 2.14.} $\underline{C}$ étant comme ci-dessus la catégorie des $A$-modules de type fini annulés par une puissance de $I$, soit $T : \underline{C}^{o} \to (\mathrm{Ab})$ un foncteur contravariant exact à gauche de $\underline{C}$ dans $(\mathrm{Ab})$, soit $(Y_i)_{i \in I}$ la famille des composantes irréductibles de $Y$ \add{$= V(I)$}, \struck{et} soit $(Y_i)_{i \in J}$ une sous-famille telle que pour tout $M \in \mathrm{Ob}\,\underline{C}$, la nullité de $T(M)$ (qui de toute façon ne dépend que de $\mathrm{supp}\, M$ en vertu de 2.4.) soit équivalente à la relation "$Y_i \not\subset \mathrm{supp}\, M$ \note{la barre de $\not\subset$ est ajoutée à l'encre} pour tout $i \in J$". Soit $S$ \add{$= (Y_i)_{i \in J}$} l'ensemble des points génériques des $Y_i$, $i \in J$. Alors pour tout $M \in \mathrm{Ob}\,\underline{C}$, on a
\[
  \mathrm{Ass}\, T(M) = \mathrm{supp}\, M \cap S ,
\]
en particulier le $A$-module $T(M)$ est sans idéaux premiers associés imm.

En effet, posant $H = \varinjlim T(A/I^n)$, on a un isomorphisme
\[
  T(M) \simeq \mathrm{Hom}(M,H)
\]
d'où
\[
  (\%) \qquad \mathrm{Ass}\, T(M) = \mathrm{supp}\, M \cap \mathrm{Ass}\, H ,
\]
et il reste à prouver la formule
\[
  \mathrm{Ass}\, H = S .
\]
\struck{Soit donc} Evidemment $\mathrm{Ass}\, H \subset \mathrm{supp}\, H \subset Y$, il reste à prouver que si $y \in Y$, alors $y \in \mathrm{Ass}\, H$ équivaut à $y \in S$. Or soit $M = A/\underline{p}$, où $\underline{p} = \underline{p}_y$ \add{(donc $\mathrm{supp}\, M = V(\underline{p}) = \bar{y}$)}, \struck{en} vertu de $(\%)$ on a $T(M) \neq 0$ (i.e. $\mathrm{Ass}\, T(M) \neq \emptyset$) sss \add{$\bar{y}$} \struck{\ill{}} rencontre \struck{pas} $\mathrm{Ass}\, H$, or par hypothèse on a $T(M) \neq 0$ sss \struck{$\mathrm{supp}\, M = V(\underline{p})$} \add{$\bar{y}$} rencontre \struck{\ill{}} $S$, ou ce qui revient au même, sss $y \in S$. \note{le $y$ porte ici une barre ajoutée à l'encre, comme aux deux lignes précédentes} Cela prouve d'abord que $\mathrm{Ass}\, H \subset S$, (prendre $y \in \mathrm{Ass}\, H$), puis $S \subset \mathrm{Ass}\, H$ puisque tout $y \in S$ doit avoir une spécialisation dans $S$, donc être dans $S$ puisque deux éléments de $S$ dont \struck{\ill{}} l'un spécialise l'autre sont identiques.

Le résultat 2.10. nous incite à traduire les propriétés des modules \add{$H$} de support $Y = V(I)$ en termes du foncteur exact à gauche correspondant $M \rightsquigarrow \mathrm{Hom}(M,H)$ sur $\underline{C}$. Ici, nous nous contentons d'expliciter le résultat suivant:

\emph{Proposition 2.15.} Avec les notations précédentes, pour que $H$ soit un $A$-module injectif, il f et s que le foncteur $T_H(M) = \mathrm{Hom}(M,H)$ sur la catégorie $\underline{C}$ des $A$-modules de type fini \struck{à supp} annulés par une puissance de $I$, soit exact à droite donc exact. En d'autres termes, 2.10. induit une équivalence entre la catégorie des foncteurs contravariants \emph{exacts} $T$ de $\underline{C}$ dans $(\mathrm{Ab})$, avec la catégorie des $A$-modules injectifs à support $\subset V(I)$ i.e. dont tout élément est annulé par une puissance de $I$.

\page{13}
Nécessité triviale, prouvons la suffisance, i.e. si $M \rightsquigarrow \mathrm{Hom}(M,H)$ est exact sur la sous-catégorie $\underline{C}$ de la catégorie des $A$-modules, alors ce foncteur est exact partout. On sait (Tohoku) qu'il suffit de le vérifier que pour tout sous-module $N$ de $M = A$, $\mathrm{Hom}(M,H) \to \mathrm{Hom}(N,H)$ est surjectif, donc on est ramené au cas où $M$, $N$ sont \struck{surjectifs} \add{de type fini}. Or un homom
\[
  u : N \longrightarrow H
\]
a comme image un sous-module de type fini de $H$, donc annulé par une puissance de $I$, i.e. \add{$u$ est} nul sur un sous-module $I^n N$ de $N$. En vertu de \emph{Krull}, $I^n N$ contient une intersection $N \cap I^m N$. Or on a une inclusion
\[
  0 \longrightarrow N/N \cap I^m M \longrightarrow M/I^m M
\]
dans $\underline{C}$, et $u$ se factorise par un homomorphisme
\[
  u' : N/N\, I^m M \longrightarrow H \; ;
\]
ce dernier par hypothèse peut se prolonger à $M/I^m M$, d'où un homomorphisme $v : M \to H$ qui prolonge $u$, cqfd.

\emph{Corollaire 2.16.} Soit $K$ un $A$-module, et soit $H^{o}_{I}(K)$ le sous-module de $K$ formé des éléments annulés par une puissance de $I$. Si $K$ est injectif, il en est de même de $H^{o}_{I}(K)$.

En effet, lorsque $M$ parcourt $\underline{C}$, on a un isomorphisme fonctoriel
\[
  \mathrm{Hom}(M, H^{o}_{I}(K)) \xrightarrow{\;\sim\;} \mathrm{Hom}(M,K) ,
\]
de sorte qu'il suffit d'appliquer le critère 2.15.

\struck{Si on rejette tout ce numéro (comme il me semble raisonnable)}

\emph{Remarques 2.17.} On peut généraliser 2.6. et 2.9., ainsi que leurs corollaires, au cas où on est sur un préschéma noethérien $X$, $\underline{C}$ désignant la catégorie des Modules cohérents sur $X$ à support $\subset Y$ ($Y = X$ dans le cas \struck{dual} de 2.6.), dans le cas d'un foncteur contravariant \add{exact à gauche} $T$ de $\underline{C}$ dans $(\mathrm{Ab})$: un tel foncteur \struck{se représente de façon} se représente (de façon essentiellement unique) comme $F \rightsquigarrow \mathrm{Hom}(F,H)$, où $H$ est un Module quasi-cohérent sur $X$ à support $\subset Y$, encore donné par \add{$H \simeq \varinjlim T(\mathcal{O}_X/I^n)$}. Pour le voir, on note que $T$ est "ind-représentable" par le critère général donné dans A. Grothendieck, Technique de descente et théorèmes d'existence en Géom Alg II, Sém Bourbaki ...., (qui s'applique ici grâce au fait que les objets de $\underline{C}$ sont noethériens), et on note que $\mathrm{Ind}(\underline{C})$ est équivalente de façon naturelle à la catégorie de tous les Modules quasi-cohérents sur $X$ ... Nous n'aurons sans doute pas besoin de ce fait dans ce Chapitre, où les résultats affines développés ici suffiront pour notre propos.

\marginal{déployer}

\note{le numéro de l'exposé Bourbaki est laissé en blanc sur la page, quatre points de suspension ; il n'est pas comblé ici}

\section{Propriétés $(P_1)$ et $(P_2)$ pour des faisceaux d'ensembles}

\note{à partir d'ici la main change : les six feuillets qui suivent sont manuscrits, d'une écriture rapide, et forment une suite numérotée par lui Définition 1, Définition 2, Proposition 3, Corollaire 4, Lemme 5. Les énoncés et les formules passent ; les phrases de liaison, souvent non}

\page{15}
Propriétés $(P_1)$ et $(P_2)$ pour des faisceaux d'ensembles.

\struck{Proposition}

\emph{Définition 1.} Soient $X$ un espace topologique, $Y$ une partie fermée de $X$, $F$ un faisceau d'ensembles sur $X$. On dit que $F$ a la propriété $(P_1)$ (resp. $(P_2)$) si le morphisme canonique
\[
  F \longrightarrow i_{*} i^{*}(F) = i_{*}(F|U)
\]
(où $i : U = X - Y \longrightarrow X$ morphisme canonique) est injectif (resp. bijectif).

(Cela signifie aussi que pour tout ouvert $V \subset X$, l'application de restriction
\[
  \Gamma(V,F) \longrightarrow \Gamma(V \cap U, F)
\]
est injective (resp. bijective)).

\marginal{N.B. regarder aussi la propriété : \uncertain{surjective} ....}

\note{suit, entre crochets et dans la colonne de gauche, un bloc que ce passage n'a pas su lire de bout en bout : il porte « Car si on prend \ill{} et \ill{} », puis « $\varinjlim \Gamma(X,F) = F_x$ », puis « injectif (resp. bijectif) »}

\ill{} \note{une phrase de transition, de cinq ou six mots, non lue : elle restreint désormais l'étude aux espaces localement noethériens dont les parties fermées \ill{}} \uncertain{noethérien loc.} \ill{} $X$ \ill{} \uncertain{[« espace noethérien spécial »]} \ill{} et \ill{} que $X$ soit un \ill{}

Définition de \struck{\ill{}} $\mathrm{Loc}_x(X)$, $x$ point d'une spécialisation. Définition de \struck{\ill{}} $\mathrm{Loc}_x(F)$, $F$ un faisceau sur un espace \ill{}

\emph{Définition 2.} Soit $F$ un faisceau \ill{} sur un espace noethérien local $X$. \ill{} On dit que $F$ a la propriété $(P_1)$, resp. $(P_2)$, si \ill{} $F$. La propriété $(P_3)$ resp. $(P_4)$ vel. \note{la fin de la définition n'est pas lue ; elle se poursuit par le cas $Y = \{x\}$ et par « si $x$ n'a \ill{} »}

\page{16}
\note{ce premier alinéa est enfermé dans une accolade à gauche}
Si $X$ est loc. noeth. quelconque, on dit que $F$ a la propriété $(P_1)$ (resp. $(P_2)$) \ill{} si pour tout $x \in X$, \uncertain{$\mathrm{Loc}_x(F)$ sur $\mathrm{Loc}_x(X)$} a la propriété $(P_1)$ resp. $(P_2)$.

\emph{Proposition 3.} Soit $X$ un espace loc. noeth. spécial, $Y$ une partie fermée de $X$, $F$ un faisceau d'ens. sur $X$. Conditions équivalentes:
\begin{itemize}
  \item (i) $F$ a la propriété $P_1$ (resp. $P_2$) vel. \ill{} $Y$
  \item (ii) \struck{\ill{}} Pour tout $x \in Y$, $F$ a la propriété $P_1$ (resp. $P_2$) en $x$.
\end{itemize}

\emph{Démonstration.} (On peut supposer $X$ affine) \struck{\ill{}} Prouvons d'abord le

\emph{Corollaire 4.} Si $F$ satisfait à la condition $P_1$ ($P_2$) rel. à $Y$, alors pour tout \struck{ouvert} \ill{} $X' \subset X$ et toute partie fermée $Y'$ de $X'$ telle que $Y' \subset X' \cap Y$, $F|X'$ satisfait $P_1$ ($P_2$) rel. à $Y'$.

\marginal{\uncertain{s'intéresser} à l'hyp. loc. noeth. sur $X$}

Soit $U' = X - Y$ \note{la ligne porte au-dessus une insertion à l'encre, non lue}. \ill{} \uncertain{On peut supposer $X' = X$}. Soit $V$ un \ill{} quelconque de $X$,

\page{17}
Pour cela \ill{}

\emph{Lemme.} Pour deux parties loc. constructibles \add{telles que $E \supset E'$} $E$, $E'$ de $X$ \struck{\ill{}} au voisinage de $x$, il f. et s. que \ill{} \struck{\ill{} $E$} \ill{} i.e.
\[
  E_{(x)} \supset E'_{(x)} .
\]

\emph{Immédiat:} En $Z = E' \cap \complement E$, on est \ill{} \ill{} que si $Z$ constructible \ill{} tel que $Z_{(x)} = \emptyset$, alors $\bar{Z} \not\ni x$, i.e. $X - Z$ est un voisinage de $x$. Immédiat.

On vient de voir que ceci est injectif. Prouvons qu'il est \emph{surjectif}. \note{« surjectif » est souligné} Pour \ill{} soit $f \in \Gamma(V_{(x)}, F_{(x)})$ \note{cette formule est ajoutée au-dessus de la ligne}. Pour tout $y \in V_{(x)}$, soit $W^y$ un voisinage ouvert de $y$, et \ill{} $f^y \in \Gamma(W^y,F)$, tel que $(f^y)_{y'} = f_{y'}$ pour $y' \in W^y_{(x)}$.

\marginal{$V_{\alpha} \subset X_{(x)}$}

Les $W^y$ recouvrent $V_{(x)}$, donc \ill{} \note{une longue rature horizontale traverse ici quatre lignes ; ce qu'elle annule n'a pas été lu}

\ill{} Pour $y' \in V_y$ voisinage de $y$ \ill{} \ill{} tel que $(f^y)_{y'} = f_{y'}$, donc \ill{} à \ill{} les \ill{} de voisinage \ill{} $V_{\alpha}$, $W^y$ \ill{} \ill{} \ill{} donc les $W^y$ recouvrent $V_{\alpha}$, \ill{} \ill{} il existe un sous-recouvrement fini des $W^y$ \ill{} \ill{}

\page{18}
On a alors le diagramme commutatif

\begin{tikzcd}
  \Gamma(V,F) \arrow[r, "\rho'"] \arrow[dr, "\rho"'] & \Gamma(V \cap W', F) \arrow[d, "\rho''"] \\
  & \Gamma(V \cap U, F)
\end{tikzcd}

Si $F$ satisfait à $P_1$ ($P_2$) rel. à $Y$, alors $\rho$ et $\rho''$ sont \struck{\ill{}} inj (bijectifs) puisque $V \cap U = (V \cap W') \cap U$, donc \struck{\ill{}} il en est \struck{\ill{}} \ill{} de même de $\rho'$ qui est inj. (bijectif) \note{la ligne est en partie annulée d'un trait ; le sens en est rétabli du contexte} ce qui prouve le corollaire.

\struck{\ill{}} \note{un intitulé biffé et souligné, d'environ quatre mots, non lu}

\emph{Lemme 5.} On a, pour tout $x \in X$, \struck{\ill{}} \marginal{et $U$ \ill{} $V$,} \ill{} \struck{\ill{}} $\mathrm{Loc}_x(X) = \mathrm{Loc}_x(X) - \{x\}$, $X_{(x)} = \mathrm{Loc}_x(X)$, \struck{$X'_x = V = \{x\}$} $V_{(x)} = V \cap X_{(x)}$, l'isomorphisme \struck{$\Gamma(X_x,F) = F_x$, $F(X'_x$}
\[
  (*) \qquad \Gamma(V_{(x)}, F_{(x)}) \;\simeq\; \varinjlim_{W \text{ voisinage ouvert de } x} \Gamma(W \cap V, F)
\]
\note{la formule $(*)$ est la seule ligne de la page écrite au net ; les lignes qui la précèdent sont un palimpseste de trois états successifs, dont la lecture ci-dessus ne retient que ce qui subsiste en dernier}

On a une lecture évidente, prouvons qu'il \note{sic} est bijectif. Deux $f$, $g \in \Gamma(W \cap V, F)$ qui deviennent \struck{\ill{}} \add{égales} \struck{\ill{}} dans $V_{(x)}$, donc ils sont \ill{} nulles dans un voisinage ouvert $E$ de $V_{(x)} \subset X$, \struck{\ill{}} je dis que $E$ contient un ens. de la forme $W' \cap V$, $W'$ ouv. \ill{} de $x$.

\page{19}
Prouvons \ill{} (ii) $\Rightarrow$ (i). La question étant locale sur $X$, on peut supposer $X$ \emph{noethérien}. On procède par récurrence noethérienne \note{« noethérienne » est souligné} \ill{} : pour toute partie fermée $Y' \subset Y$, $F$ \ill{} la propriété $P_1$ ($P_2$) rel. à $Y'$, on \ill{} supposer \ill{} que \ill{} pour tout $Y'$ fermé $\subset Y$, $Y' \neq Y$, \marginal{rel. à $Y'$} \ill{} $F$ a déjà la propriété en question \ill{}, et borner : prouver \ill{} $F$ ait la propriété $P_1$ ($P_2$) rel. à $Y$ \ill{} \struck{\ill{}} Si $Y = \emptyset$, c'est trivial, \struck{\ill{}} sinon \struck{\ill{}} \ill{} soit \ill{} une \emph{partie} de $Y$. \ill{} \ill{} \ill{} il faut prouver que
\[
  \Gamma(X,F) \longrightarrow \Gamma(U,F) \qquad (U = X - Y)
\]
est injectif (bijectif). Soient $f$, $g \in \Gamma(X,F)$ ayant \ill{} même image dans $\Gamma(U,F)$, prouvons qu'ils sont \struck{identiques} \ill{} \struck{\ill{}} \ill{} \struck{il \ill{} $x \in Y$,} $f_x = g_x$ si \ill{} $\bar{x} \neq Y$ \marginal{et \ill{} l'hyp. $P_1$ \ill{} $x$} local. Tout d'abord, on \struck{\ill{}} \ill{} ils sont identiques en $x$, donc un voisinage de $x$, et \ill{} \ill{} par l'hyp. de récurrence \ill{} \ill{} \ill{} (et l'hyp. $P_2$ \ill{} $x \in Y$)

\note{la page 19 est le feuillet le plus travaillé du lot : trois longues ratures horizontales, deux blocs cernés, un mot biffé dans la marge gauche. La structure de l'argument — récurrence noethérienne sur les fermés de $Y$, réduction à l'injectivité (bijectivité) de $\Gamma(X,F) \to \Gamma(U,F)$ — passe ; le détail des phrases, souvent non}

\page{20}
\note{une insertion à l'encre en tête de page, cernée et en partie biffée, n'a pas été lue}
\ill{} éléments $W_i$ de $X$ recouvrant $V_{(x)}$ \ill{} et des $f_i \in \Gamma(W_i,F)$, tels que $y \in W_{i,x}$ \struck{\ill{}} $\Rightarrow$ $f_{i,y} = f_y$.

Soit \struck{\ill{}} $W_{ij}$ l'ens. des $y \in W_i \cap W_j = W_{ij}$ tels que $f_{i,y} = h_y$. \note{sic : le second membre est écrit $h$ et non $f_j$} C'est un \struck{\ill{}} \add{\ill{}} constructible \ill{}
\[
  W_{ij} \cap V_{(x)} = W_{ij,x}
\]
Donc \ill{} \ill{} pour $i$ et $j$, \ill{} relativement à $X$ \ill{} \ill{} de \struck{\ill{}}, on peut supposer que
\[
  W'_{ij} = W_{ij} .
\]
Alors \ill{} \ill{} \ill{} \ill{} une vraie section. OK ....

\marginal{généralisation \ill{} plus généraux}

Prouvons maintenant (i) $\Rightarrow$ (ii) : en effet, soit $x \in Y$, \struck{\ill{}} \ill{} \ill{} \emph{seul} \ill{} corollaire 4, $F$ satisfait à $P_1$ ($P_2$) rel. à $\bar{x}$, donc \add{avec $V = X - \bar{x}$} \ill{} \ill{} \ill{} pour tout voisinage ouvert $W$ de $x$
\[
  \Gamma(W,F) \longrightarrow \Gamma(W \cap V, F)
\]
est injective (bij.), donc
\[
  \varinjlim_W \Gamma(W,F) \longrightarrow \Gamma(W \cap V, F)
\]
est injective (bij.) et en vertu \ill{} s'identifie à :
\[
  \Gamma(X_{(x)}, F_{(x)}) \longrightarrow \Gamma(X'_{(x)}, F_{(x)})
\]
\ill{} \uncertain{en appliquant deux fois le Lemme 5}. OK.

\end{document}
